\paragraph{Theorem F.2.29 (Proved: a complete finite integral survivor test).}
Take any rational pair in F.2.20 and any of its integral,
primitive-cusp rank-six extensions supplied by F.2.22.
Write $M_1,M_2$ for geometric transport on fibre homology
$\Lambda=\mathbb Z^6$, in the peripheral convention
$M_0M_1M_2=I$.
First require the preliminary conditions
\[
 \Lambda_\Gamma=\mathbb Z,\qquad
 \ell(\Lambda^{M_i})=\mathbb Z,\qquad h_1=5,\quad h_2=4 .
 \tag{F.HC1}
\]
Here $\ell$ is the primitive global invariant covector, and
$h_i$ is the integral fixed-vector/fixed-covector pairing
on $\ker\ell$ from F.2.25.
Failure of any condition excludes the datum.
If they hold, use the unique normalized affine filling of
F.2.28, with twists $(u_1,-u_2)$ and zero cusp clutching.
The following six finite integer-matrix tests are necessary
and sufficient for this filling to have $H_2=H_3=0$.

Let
\[
 T_i=M_i^{-t},\quad T_0=(T_1T_2)^{-1},\qquad
 G_q=\operatorname{coker}
 [\,\bigwedge^qT_1-I\mid\bigwedge^qT_2-I\,]
 \quad(1\le q\le4).
 \tag{F.HC2}
\]
All $G_q$ are finite.  Let $P_q$ be the integral parabolic
cohomology group of the exterior local system
$\bigwedge^q\mathcal V$, for $q=2,4$.
Let $t_q\in G_{q-1}$, $2\le q\le5$, be the natural global
coinvariant image of the unique nonzero torsion suspension
class at the order-four cap, as specified below.
Let $c_4\in G_3$ be the natural image of the unique nonzero
torsion class in
$\operatorname{coker}(\bigwedge^3T_0-I)$.
Then the six tests are
\[
\begin{array}{c|ccc}
 &\text{first test}&\text{second test}&\text{third test}\\ \hline
 H^3=0&G_1=\langle t_2\rangle&P_2=0&t_3\ne0\\
 H^5=0&G_3=\langle t_4,c_4\rangle&P_4=0&t_5\ne0 .
\end{array}
\tag{F.HC3}
\]
Every group, class, and equality here is computed by
finitely many integer kernels, Smith forms, and exterior
minors.  This is a decision criterion for every rational
parameter and every integral lattice in the complete
normal form; it makes no bounded-height assertion about
the rational parameter set.

\begin{proof}
We first establish the uniform local calculations needed
for the criterion, so the test does not assume that a
general finite matrix is globally conjugate to the example
in F.2.6.
The fixed vectors $u_i$ split the homology lattice as
$\mathbb Zu_i\oplus H_i$, where $H_i=\ker\ell$.
Write $D_i=M_i|_{H_i}$.
The maximal pairing $h_i=m_i$ gives
\[
 \operatorname{Norm}_{m_i}(D_i)=\lambda_i\varphi_i ,
 \tag{F.HC4}
\]
with both factors primitive.

At five, work over $\mathbb Z_5$.
Modulo five the norm is $(D_1-I)^4$ and is nonzero.
Choose $q$ with $\operatorname{Norm}_5q$ primitive.
The five vectors $q,D_1q,\ldots,D_1^4q$ are independent
modulo five: otherwise the nilpotent Jordan block generated
by $q$ would have length less than five.
They are therefore a $\mathbb Z_5$ basis.
Thus $D_1$ is the regular five-cycle permutation locally.

At four, work over $\mathbb Z_2$.
The norm is $(D_2-I)^3$ modulo two and is nonzero.
The vectors $q,D_2q,D_2^2q,D_2^3q$ similarly generate
a primitive regular four-cycle submodule $Q$ of rank four.
The quotient is rank one with $D_2=-1$, by the
characteristic polynomial.
For a lift $z$ of its generator, put $a=(D_2+I)z\in Q$.
Since $D_2^4=I$,
\[
 (D_2^3-D_2^2+D_2-I)a=0.
\]
On the regular four-cycle module the kernel of the
displayed alternating norm is exactly $(D_2+I)Q$,
integrally.  One can see this by writing its four
coefficients; equivalently $D_2+I$ has three nonzero
Smith factors, all one, and the indicated rational image.
Subtracting a suitable element of $Q$ from $z$ makes
$D_2z=-z$.
Consequently
\[
 M_1\sim 1\oplus\operatorname{perm}_5\quad\text{over }\mathbb Z_5,
 \qquad
 M_2\sim1\oplus\operatorname{perm}_4\oplus(-1)
 \quad\text{over }\mathbb Z_2 .
 \tag{F.HC5}
\]
These local comparisons apply to cap cohomology, not just
to abstract representation lattices.
Indeed the integral lattice of intertwiners is the kernel
of an integer matrix and its tensor with $\mathbb Z_p$
is the lattice of local intertwiners.
Approximate a local invertible intertwiner by an integral
one whose determinant is a $p$-adic unit.
Keeping the $u_i$ summand fixed produces an equivariant
torus isogeny of degree prime to $p$, preserving the
normalized translation $\pm u_i/m_i$.
The resulting cap and boundary coverings induce
isomorphisms on $\mathbb Z_p$ cohomology: the vertical
torus maps are cohomology isomorphisms over $\mathbb Z_p$,
so the mapping-torus Wang sequences give the assertion.
At primes not dividing $m_i$, transfer already makes
finite-cover pullback an isomorphism onto invariants.

Here are the full required integral results for these
two monomial models.
The ranks of their invariant exterior lattices in
degrees zero through six are
$(1,2,3,4,3,2,1)$.
The order-five cap has no cohomology torsion in degrees
two through five.
The order-four cap has exactly $\mathbb Z/2$ in each
of those four degrees.
This follows from the Wang sequence for the mapping
torus of the vertical $D_i$:
its torsion in degree $q$ is
\[
 \operatorname{Tor}\operatorname{coker}
       (\bigwedge^{q-1}D_i^t-I).
 \tag{F.HC6}
\]
In particular it introduces no additional torsion primes.

Integral pullback from either finite cap is onto the
invariant lattice in degrees two through five.
For degree two this is precisely the $m_i/h_i=1$
suspension calculation of F.2.25.
For completeness, the determinants of the pairings between
the invariant lattices in complementary degrees
$(1,5),(2,4),(3,3)$ are respectively
$m_i,m_i^3,m_i^4$ in absolute value.
The degree-one cover image has index $m_i$.
The free downstairs complementary pairings are unimodular,
and pullback multiplies each pairing by the covering
degree $m_i$.
The middle pairing first gives degree-three index one.
The degree-two index one then gives degree-four index one;
the degree-one index $m_i$ gives degree-five index one.
These determinant identities and all the Wang Smith
forms are verified by the fixed, exhaustive computation
in \texttt{fourfolds\_local\_caps.py}, for both orders and
all exterior degrees zero through six.
Thus no unexamined finite matrix is being inferred from
an example.

We next identify the boundary maps.
Choose any unimodular matrix $P_2=[u_2,H_2]$ realizing
the splitting at four, and put
$J=P_2^{-t}|_{\mathbb Z^5}$ for the inclusion of vertical
cohomology.
Choose the unique nonzero torsion representative
\[
 w_q\in\operatorname{coker}
            (\bigwedge^{q-1}D_2^t-I),\qquad
 t_q=[(\bigwedge^{q-1}J)w_q]\in G_{q-1}.
 \tag{F.HC7}
\]
A Smith transformation supplies an integral $w_q$.
Its norm is exactly zero, since it lies rationally in
the image of $\bigwedge^{q-1}D_2^t-I$.
The finite boundary coordinates of F.2.6 satisfy
$\theta=a+4k$, $h=k$, hence $a=\theta-4h$.
The pullback of the suspension cocycle has $\theta$
component $w_q$ and $h$ component
$-\operatorname{Norm}_4(w_q)=0$.
This holds in every exterior degree, including $q=5$.
It is an integral cellular assertion for the monomial
model, where the torus cellular differentials vanish;
the prime-to-two isogeny comparison transports it to
every eligible cap.
The base connecting map therefore gives exactly the
class $t_q$ in (F.HC7).
Different representative choices change it by a global
monodromy relation.  Any sign change is immaterial for
an order-two class.

There is one further boundary map, from the free
degree-four cusp kernel.
The primitive cusp is integrally conjugate to
$\left(\begin{smallmatrix}I&I\\0&I\end{smallmatrix}\right)$.
Its degree-three exterior difference has ten Smith
factors one and one Smith factor two; its cokernel is
$\mathbb Z^9\oplus\mathbb Z/2$.
The other exterior differences in degrees one, two,
four and five have torsion-free cokernels.
The exact matrices and a torsion representative are
saved by \texttt{fourfolds\_cusp\_boundary.py}.
Let $C$ be the smooth cusp cap and $Y=\partial C$.
The boundary Wang sequence gives
$\operatorname{Tor}H^4(Y)=\mathbb Z/2$.
In the Poincare--Lefschetz exact sequence its image in
\[
 H^5(C,Y)\cong H_3(C)=\mathbb Z^9
\]
is zero, so this torsion class lifts to $H^4(C)$.
Its restriction to an ordinary fibre is zero.
By E.1.3 the kernel of integral specialization is
a primitive copy $K=\mathbb Z$ in degree four.
The lift is in $K$, and the image of this cyclic group
in $H^4(Y)$ contains the nonzero torsion class.
A cyclic image containing nonzero torsion cannot
also have a nonzero free part.
Its primitive generator therefore restricts to the
unique nonzero torsion class of $H^4(Y)$.
Consequently the cusp contribution to $d_2$ is
exactly $c_4$ as defined above.  In particular this
column must not be discarded merely because it
vanishes in the first examples.

It remains to prove that these are all the tests.
For the irreducible odd-dimensional quotient $W$
of $\mathcal V$, invariant alternating forms and
bivectors vanish: a nonzero one would give an
isomorphism $W\cong W^*$ with an odd-dimensional
nondegenerate alternating pairing.
The exact exterior sequences for
$0\to\mathbb Q\to\mathcal V_{\mathbb Q}\to W\to0$,
together with $\det W=1$, give
\[
 H^0(B,j_*\bigwedge^q\mathcal V)=0
       \quad(2\le q\le5),\qquad
 G_q\otimes\mathbb Q=0\quad(1\le q\le4).
 \tag{F.HC8}
\]
For the endpoint assertions, the extension cannot split:
on $\mathbb Q\oplus W$ a square-zero nilpotent has rank at
most two, whereas the stipulated rank is three.
An invariant covector of $\mathcal V_{\mathbb Q}$ would
split this extension, since $W$ has no trivial quotient.
Thus $(\mathcal V_{\mathbb Q}^*)^\Gamma=0$, which gives
both $H^0(\bigwedge^5\mathcal V_{\mathbb Q})=0$ and
$G_1\otimes\mathbb Q=0$.
Here and below the first assertion is integral as well,
since the coefficient lattices are torsion free.
The local invariant ranks in exterior degrees two
and four are $3,3,9$ at the three punctures.
The Euler formula for the rank-fifteen local system
therefore gives $P_2\otimes\mathbb Q=P_4\otimes\mathbb Q=0$.
In degree one it similarly gives rational parabolic
cohomology zero.
The only global degree-one section is the invariant
line.  Simple connectivity and the universal
coefficient theorem make its free transgression
onto $H^2(B,\mathbb Z)$; equivalently the rational
second Betti number is zero, as in F.2.6.
Thus the normalized filling has rational $b_2=b_3=0$.

All specializations in exterior degrees two through
five are onto, by the local calculations and E.1.3.
Their only point-supported kernels are
\[
\begin{array}{c|cccc}
q&2&3&4&5\\ \hline
\text{kernel}&\mathbb Z/2&\mathbb Z/2&
                 \mathbb Z\oplus\mathbb Z/2&\mathbb Z/2 .
\end{array}
\tag{F.HC9}
\]
The free summand in degree four is at the cusp; the
others are at four.
Consequently $H^1(B,R^qg_*\mathbb Z)=P_q$ in the
degrees used here, while
$H^2(B,R^qg_*\mathbb Z)=G_q$.
The latter equality also holds in degree one:
changing a sheaf at isolated points changes no
degree-two cohomology.  The equality with coinvariants
follows from compactly supported top cohomology on
the punctured sphere and the sequence $j_!\to j_*$.

The base has real dimension two.
Thus the only possible nonzero differential in the
integral Leray spectral sequence is
$d_2:E_2^{0,q}\to E_2^{2,q-1}$.
The complete associated graded groups in total
degrees three and five are therefore
\[
\begin{split}
 \operatorname{gr}H^3&:
   G_1/\langle t_2\rangle,\quad P_2,\quad
   \ker(\mathbb Z/2\xrightarrow{\,t_3\,}G_2),\\
 \operatorname{gr}H^5&:
   G_3/\langle t_4,c_4\rangle,\quad P_4,\quad
   \ker(\mathbb Z/2\xrightarrow{\,t_5\,}G_4).
\end{split}
\tag{F.HC10}
\]
These are all finite.  A finite filtered group
vanishes exactly when every associated graded group
vanishes, regardless of extension classes.
Since $\pi_1=1$ and rational $b_2=b_3=0$, the universal
coefficient theorem gives
$H^3=\operatorname{Ext}(H_2,\mathbb Z)$, so
$H_2=0$ exactly when $H^3=0$.
Integral Poincare duality gives $H^5\cong H_3$.
This proves (F.HC3).

For reproducibility, no sheaf calculation is left
as a black box in the integer test.
Given $A=\bigwedge^qT_1$, $B=\bigwedge^qT_2$ on a
free lattice $L$, form the saturated integer kernel
\[
 K_q=\ker[\,A-I\mid A(B-I)\mid -(AB-I)\,]\subset L^3.
 \tag{F.HC11}
\]
Its image under
$(x,y,z)\mapsto((A-I)x,(B-I)y)$ is precisely the
integral parabolic cocycle lattice.
Quotient by the columns of
$\left(\begin{smallmatrix}A-I\\B-I\end{smallmatrix}\right)$
to obtain $P_q$.
This uses the actual integral images of peripheral
differences, not their saturations.
Smith transformations also supply every $w_q$, the
cusp order-two representative, and the tests for
their images and generated subgroups in $G_q$.

The program \texttt{fourfolds\_survivor\_test.py}
implements exactly these operations.
Its executed test set is precisely $v=0,1,2$ from
F.2.16.  For $v=0,1$ all six tests pass;
$G_1,\ldots,G_4$ are each $\mathbb Z/2$,
$P_2=P_4=0$, all four finite-cap images are nonzero,
and the cusp image is zero.
For $v=2$ they are each $\mathbb Z/6$,
$P_2=P_4=\mathbb Z/3$, and the first, second,
fourth and fifth tests fail.  The computation thus
checks both accepted and rejected data.

We finish by making the convention dictionary explicit.
The imported family matrices called $A,B$ in that
program are geometric cohomology pullbacks, so the
geometric homology inputs are $(A^t,B^t)$.
Their genuine left cohomology transport in (F.HC2)
is $(A^{-1},B^{-1})$, which is what the code uses.
For comparison with computations written using
$(A,B)$, simultaneous reversal of the two peripheral
generators sends
\[
 g_1\mapsto g_1^{-1},\qquad
 g_2\mapsto g_2^{-1},\qquad
 g_0=(g_1g_2)^{-1}\mapsto
     g_2g_1=g_1^{-1}(g_1g_2)g_1 .
 \tag{F.HC12}
\]
This is an orientation-reversing marking of the
pair of pants fixing its three ends: each peripheral
conjugacy class is inverted.
Thus the new cusp is conjugate to the inverse old
cusp, not an independently chosen third matrix.
Inversion preserves invariant and coinvariant
lattices, and the indicated conjugating generator
acts trivially on global coinvariants.
All boundary classes in (F.HC3) have order two, so
reversal signs do not change the tests.
The implementation uses (F.HC2) directly and hence
does not depend on this comparison for its validity.
\end{proof}

\paragraph{Corollary F.2.30 (Proved: constructive classification of all survivors).}
For the marked Kuhn cap and affine translation clutchings
used throughout F.2, the following is a complete constructive
classification of the data with
$\pi_1=1$ and $H_2=H_3=0$, up to integral conjugacy of their
linear tuples and oriented diffeomorphism of their fillings.

There are no data with three genuine finite fillings.
With two fillings their actual affine orders are five and
four.  Order the order-five point first.
The complete, unbounded parameter set is the rational
projective normal form of F.2.20 together with its
exceptional rational affine line.
For each parameter, run the terminating order,
conductor, stable-lattice, and primitive-extension algorithm
of F.2.22.  Discard precisely the integral classes failing
(F.HC1) or one of (F.HC3).
The classes retained by this procedure are exactly all
surviving integral linear tuples.
For each retained class all simply connected affine
twists are exactly the choices satisfying (F.T12);
their fillings are all oriented diffeomorphic to its
normalized filling.
If ordered labels with the order-four point first are
desired, interchange the finite labels and recompute
$M_0=(M_1M_2)^{-1}$.

Every retained datum has
\[
 \pi_1=1,\qquad
 H_j=\begin{cases}
  \mathbb Z,&j=0,8,\\
  \mathbb Z^4,&j=4,\\
  0,&\text{otherwise},
 \end{cases}
 \tag{F.HC13}
\]
and its oriented smooth manifold is
$3\mathbb{HP}^2\#(-\mathbb{HP}^2)$.
It has bounded-cusp holomorphic periods in the literal
sense of F.3, by F.3.11.
The classification here concerns the survivor subproblem.
It does not classify the nonsurviving tuples or compute
their full integral homology, as also requested in F.2.

\begin{proof}
The exclusions F.2.7, F.2.14, F.2.25, and F.2.27
leave exactly the stated minimal orders and rational
quotient problem.
The rational normal form is injective on conjugacy
classes, including the exceptional line.
The finite interval and extension-coset enumeration
in F.2.22 gives every integral class over each rational
point, with all bounds proved for that input.
No distinct rational points can become integrally
conjugate, since an integral conjugacy induces a
rational quotient conjugacy.
Theorem F.2.28 and the necessary-and-sufficient
criterion F.2.29 prove both directions of the
retention rule and account for every affine twist.
This is also an exhaustive enumeration if desired:
enumerate primitive projective integer triples with
first nonzero coordinate positive by increasing height,
and enumerate reduced rational exceptional parameters
by increasing numerator and denominator height.
Every finite stage has finitely many inputs, each
terminates by F.2.22, and every rational input occurs.
No finite cutoff is imposed or needed.

Hurewicz makes the retained simply connected manifolds
three-connected.  The only singular contribution to
Euler number is the Kuhn cusp, whose value is six
by E.1.3; torus bundles and the free finite caps have
Euler number zero.
Integral Poincare duality and the universal coefficient
theorem then give (F.HC13), including torsion freeness.

For the smooth classification, F.2.15 supplies a
compatible complex structure at each finite cap, and
the cusp already has its complex model.
Choose their common fibre orientation to be the one
induced by the cusp.
Commuting complex structures at the finite caps can
be chosen in that component: their two-dimensional
fixed real summands allow the orientation to be
reversed when necessary.
The fibre of the bundle of orientation-compatible
complex structures retracts to $\mathbb{CP}^3$.
Its connectedness and simple connectivity extend
the prescribed boundary structures across the
two-dimensional pair of pants, exactly as in F.2.5.
A compatible horizontal distribution gives an
almost complex structure on the total manifold.
Theorem F.1.2 therefore gives signature two,
intersection lattice $I_{3,1}$, and $c_2^2=2$.
Since $H^2=0$, $p_1/2=-c_2$ is primitive,
characteristic, and has square two.
The complete Wall--Nikulin--Nagy argument in F.4.1
depends only on these almost complex characteristic
numbers and three-connectivity.
It applies without alteration and gives the claimed
exact oriented diffeomorphism.
Finally F.3.11 applies to the resulting determinant-one
monodromy and primitive cusp.
It asserts the requested punctured-base periods with
bounded correction, not an additional unproved
simultaneous good extension at all finite points.
\end{proof}
